% Options for packages loaded elsewhere
\PassOptionsToPackage{unicode}{hyperref}
\PassOptionsToPackage{hyphens}{url}
\documentclass[
]{article}
\usepackage{xcolor}
\usepackage[margin=1in]{geometry}
\usepackage{amsmath,amssymb}
\setcounter{secnumdepth}{-\maxdimen} % remove section numbering
\usepackage{iftex}
\ifPDFTeX
  \usepackage[T1]{fontenc}
  \usepackage[utf8]{inputenc}
  \usepackage{textcomp} % provide euro and other symbols
\else % if luatex or xetex
  \usepackage{unicode-math} % this also loads fontspec
  \defaultfontfeatures{Scale=MatchLowercase}
  \defaultfontfeatures[\rmfamily]{Ligatures=TeX,Scale=1}
\fi
\usepackage{lmodern}
\ifPDFTeX\else
  % xetex/luatex font selection
\fi
% Use upquote if available, for straight quotes in verbatim environments
\IfFileExists{upquote.sty}{\usepackage{upquote}}{}
\IfFileExists{microtype.sty}{% use microtype if available
  \usepackage[]{microtype}
  \UseMicrotypeSet[protrusion]{basicmath} % disable protrusion for tt fonts
}{}
\makeatletter
\@ifundefined{KOMAClassName}{% if non-KOMA class
  \IfFileExists{parskip.sty}{%
    \usepackage{parskip}
  }{% else
    \setlength{\parindent}{0pt}
    \setlength{\parskip}{6pt plus 2pt minus 1pt}}
}{% if KOMA class
  \KOMAoptions{parskip=half}}
\makeatother
\usepackage{color}
\usepackage{fancyvrb}
\newcommand{\VerbBar}{|}
\newcommand{\VERB}{\Verb[commandchars=\\\{\}]}
\DefineVerbatimEnvironment{Highlighting}{Verbatim}{commandchars=\\\{\}}
% Add ',fontsize=\small' for more characters per line
\usepackage{framed}
\definecolor{shadecolor}{RGB}{248,248,248}
\newenvironment{Shaded}{\begin{snugshade}}{\end{snugshade}}
\newcommand{\AlertTok}[1]{\textcolor[rgb]{0.94,0.16,0.16}{#1}}
\newcommand{\AnnotationTok}[1]{\textcolor[rgb]{0.56,0.35,0.01}{\textbf{\textit{#1}}}}
\newcommand{\AttributeTok}[1]{\textcolor[rgb]{0.13,0.29,0.53}{#1}}
\newcommand{\BaseNTok}[1]{\textcolor[rgb]{0.00,0.00,0.81}{#1}}
\newcommand{\BuiltInTok}[1]{#1}
\newcommand{\CharTok}[1]{\textcolor[rgb]{0.31,0.60,0.02}{#1}}
\newcommand{\CommentTok}[1]{\textcolor[rgb]{0.56,0.35,0.01}{\textit{#1}}}
\newcommand{\CommentVarTok}[1]{\textcolor[rgb]{0.56,0.35,0.01}{\textbf{\textit{#1}}}}
\newcommand{\ConstantTok}[1]{\textcolor[rgb]{0.56,0.35,0.01}{#1}}
\newcommand{\ControlFlowTok}[1]{\textcolor[rgb]{0.13,0.29,0.53}{\textbf{#1}}}
\newcommand{\DataTypeTok}[1]{\textcolor[rgb]{0.13,0.29,0.53}{#1}}
\newcommand{\DecValTok}[1]{\textcolor[rgb]{0.00,0.00,0.81}{#1}}
\newcommand{\DocumentationTok}[1]{\textcolor[rgb]{0.56,0.35,0.01}{\textbf{\textit{#1}}}}
\newcommand{\ErrorTok}[1]{\textcolor[rgb]{0.64,0.00,0.00}{\textbf{#1}}}
\newcommand{\ExtensionTok}[1]{#1}
\newcommand{\FloatTok}[1]{\textcolor[rgb]{0.00,0.00,0.81}{#1}}
\newcommand{\FunctionTok}[1]{\textcolor[rgb]{0.13,0.29,0.53}{\textbf{#1}}}
\newcommand{\ImportTok}[1]{#1}
\newcommand{\InformationTok}[1]{\textcolor[rgb]{0.56,0.35,0.01}{\textbf{\textit{#1}}}}
\newcommand{\KeywordTok}[1]{\textcolor[rgb]{0.13,0.29,0.53}{\textbf{#1}}}
\newcommand{\NormalTok}[1]{#1}
\newcommand{\OperatorTok}[1]{\textcolor[rgb]{0.81,0.36,0.00}{\textbf{#1}}}
\newcommand{\OtherTok}[1]{\textcolor[rgb]{0.56,0.35,0.01}{#1}}
\newcommand{\PreprocessorTok}[1]{\textcolor[rgb]{0.56,0.35,0.01}{\textit{#1}}}
\newcommand{\RegionMarkerTok}[1]{#1}
\newcommand{\SpecialCharTok}[1]{\textcolor[rgb]{0.81,0.36,0.00}{\textbf{#1}}}
\newcommand{\SpecialStringTok}[1]{\textcolor[rgb]{0.31,0.60,0.02}{#1}}
\newcommand{\StringTok}[1]{\textcolor[rgb]{0.31,0.60,0.02}{#1}}
\newcommand{\VariableTok}[1]{\textcolor[rgb]{0.00,0.00,0.00}{#1}}
\newcommand{\VerbatimStringTok}[1]{\textcolor[rgb]{0.31,0.60,0.02}{#1}}
\newcommand{\WarningTok}[1]{\textcolor[rgb]{0.56,0.35,0.01}{\textbf{\textit{#1}}}}
\usepackage{graphicx}
\makeatletter
\newsavebox\pandoc@box
\newcommand*\pandocbounded[1]{% scales image to fit in text height/width
  \sbox\pandoc@box{#1}%
  \Gscale@div\@tempa{\textheight}{\dimexpr\ht\pandoc@box+\dp\pandoc@box\relax}%
  \Gscale@div\@tempb{\linewidth}{\wd\pandoc@box}%
  \ifdim\@tempb\p@<\@tempa\p@\let\@tempa\@tempb\fi% select the smaller of both
  \ifdim\@tempa\p@<\p@\scalebox{\@tempa}{\usebox\pandoc@box}%
  \else\usebox{\pandoc@box}%
  \fi%
}
% Set default figure placement to htbp
\def\fps@figure{htbp}
\makeatother
\setlength{\emergencystretch}{3em} % prevent overfull lines
\providecommand{\tightlist}{%
  \setlength{\itemsep}{0pt}\setlength{\parskip}{0pt}}
\usepackage{bookmark}
\IfFileExists{xurl.sty}{\usepackage{xurl}}{} % add URL line breaks if available
\urlstyle{same}
% fallback for those not using the hyperref driver hyperxmp:
\makeatletter
\@ifundefined{xmpquote}{\newcommand{\xmpquote}[1]{#1}}{}
\makeatother
\hypersetup{
  pdftitle={RPractical\_Gania},
  pdfauthor={Jerry May Gania},
  hidelinks,
  pdfcreator={LaTeX via pandoc}}

\title{RPractical\_Gania}
\author{Jerry May Gania}
\date{2026-11-09}

\begin{document}
\maketitle

\subsubsection{Part A - Basic
Operations}\label{part-a---basic-operations}

\begin{Shaded}
\begin{Highlighting}[]
\CommentTok{\#A1. Create v as the sequence 5, 10, 15, ..., 50 using seq(). Print its length, sum and mean.}

\NormalTok{v }\OtherTok{\textless{}{-}} \FunctionTok{seq}\NormalTok{(}\DecValTok{5}\NormalTok{, }\DecValTok{50}\NormalTok{, }\AttributeTok{by =} \DecValTok{5}\NormalTok{)}

\NormalTok{v}
\end{Highlighting}
\end{Shaded}

\begin{verbatim}
##  [1]  5 10 15 20 25 30 35 40 45 50
\end{verbatim}

\begin{Shaded}
\begin{Highlighting}[]
\FunctionTok{length}\NormalTok{(v)}
\end{Highlighting}
\end{Shaded}

\begin{verbatim}
## [1] 10
\end{verbatim}

\begin{Shaded}
\begin{Highlighting}[]
\FunctionTok{sum}\NormalTok{(v)}
\end{Highlighting}
\end{Shaded}

\begin{verbatim}
## [1] 275
\end{verbatim}

\begin{Shaded}
\begin{Highlighting}[]
\FunctionTok{mean}\NormalTok{(v)}
\end{Highlighting}
\end{Shaded}

\begin{verbatim}
## [1] 27.5
\end{verbatim}

\begin{Shaded}
\begin{Highlighting}[]
\CommentTok{\#A2. Create w that repeats the vector c(1, 2, 3) four times and z that repeats each element of c(1, 2, 3) four times. Print both.}

\NormalTok{w }\OtherTok{\textless{}{-}} \FunctionTok{rep}\NormalTok{(}\FunctionTok{c}\NormalTok{(}\DecValTok{1}\NormalTok{, }\DecValTok{2}\NormalTok{, }\DecValTok{3}\NormalTok{), }\AttributeTok{times =} \DecValTok{4}\NormalTok{)}
\NormalTok{z }\OtherTok{\textless{}{-}} \FunctionTok{rep}\NormalTok{(}\FunctionTok{c}\NormalTok{(}\DecValTok{1}\NormalTok{, }\DecValTok{2}\NormalTok{, }\DecValTok{3}\NormalTok{), }\AttributeTok{each =} \DecValTok{4}\NormalTok{)}

\NormalTok{w}
\end{Highlighting}
\end{Shaded}

\begin{verbatim}
##  [1] 1 2 3 1 2 3 1 2 3 1 2 3
\end{verbatim}

\begin{Shaded}
\begin{Highlighting}[]
\NormalTok{z}
\end{Highlighting}
\end{Shaded}

\begin{verbatim}
##  [1] 1 1 1 1 2 2 2 2 3 3 3 3
\end{verbatim}

\begin{Shaded}
\begin{Highlighting}[]
\CommentTok{\#A3. Compute and print: the remainder of 17 divided by 5, the integer division of 17 by 5, 2 to the power of 10, the square root of 144, and pi rounded to 3 decimals.}

\DecValTok{17} \SpecialCharTok{\%\%} \DecValTok{5}
\end{Highlighting}
\end{Shaded}

\begin{verbatim}
## [1] 2
\end{verbatim}

\begin{Shaded}
\begin{Highlighting}[]
\DecValTok{17} \SpecialCharTok{\%/\%} \DecValTok{5}
\end{Highlighting}
\end{Shaded}

\begin{verbatim}
## [1] 3
\end{verbatim}

\begin{Shaded}
\begin{Highlighting}[]
\DecValTok{2} \SpecialCharTok{\^{}} \DecValTok{10}
\end{Highlighting}
\end{Shaded}

\begin{verbatim}
## [1] 1024
\end{verbatim}

\begin{Shaded}
\begin{Highlighting}[]
\FunctionTok{sqrt}\NormalTok{(}\DecValTok{144}\NormalTok{)}
\end{Highlighting}
\end{Shaded}

\begin{verbatim}
## [1] 12
\end{verbatim}

\begin{Shaded}
\begin{Highlighting}[]
\FunctionTok{round}\NormalTok{(pi, }\DecValTok{3}\NormalTok{)}
\end{Highlighting}
\end{Shaded}

\begin{verbatim}
## [1] 3.142
\end{verbatim}

\begin{Shaded}
\begin{Highlighting}[]
\CommentTok{\#A4.Given a \textless{}{-} c(2, 4, 6, 8) and b \textless{}{-} c(1, 3, 5, 7)}

\NormalTok{a }\OtherTok{\textless{}{-}} \FunctionTok{c}\NormalTok{(}\DecValTok{2}\NormalTok{, }\DecValTok{4}\NormalTok{, }\DecValTok{6}\NormalTok{, }\DecValTok{8}\NormalTok{)}
\NormalTok{b }\OtherTok{\textless{}{-}} \FunctionTok{c}\NormalTok{(}\DecValTok{1}\NormalTok{, }\DecValTok{3}\NormalTok{, }\DecValTok{5}\NormalTok{, }\DecValTok{7}\NormalTok{)}
\NormalTok{a }\SpecialCharTok{+}\NormalTok{ b}
\end{Highlighting}
\end{Shaded}

\begin{verbatim}
## [1]  3  7 11 15
\end{verbatim}

\begin{Shaded}
\begin{Highlighting}[]
\NormalTok{a }\SpecialCharTok{*}\NormalTok{ b}
\end{Highlighting}
\end{Shaded}

\begin{verbatim}
## [1]  2 12 30 56
\end{verbatim}

\begin{Shaded}
\begin{Highlighting}[]
\NormalTok{a }\SpecialCharTok{\textgreater{}}\NormalTok{ b}
\end{Highlighting}
\end{Shaded}

\begin{verbatim}
## [1] TRUE TRUE TRUE TRUE
\end{verbatim}

\begin{Shaded}
\begin{Highlighting}[]
\FunctionTok{sum}\NormalTok{(a }\SpecialCharTok{*}\NormalTok{ b)}
\end{Highlighting}
\end{Shaded}

\begin{verbatim}
## [1] 100
\end{verbatim}

\begin{Shaded}
\begin{Highlighting}[]
\CommentTok{\#A5. Use paste() to produce "Student{-}1, "Student{-}2, "Student{-}3". Then identify the datatype 42, 42L, 41i, and TRUE. What are its output?}

\FunctionTok{paste}\NormalTok{(}\StringTok{"Student"}\NormalTok{, }\DecValTok{1}\SpecialCharTok{:}\DecValTok{3}\NormalTok{, }\AttributeTok{sep =} \StringTok{"{-}"}\NormalTok{)}
\end{Highlighting}
\end{Shaded}

\begin{verbatim}
## [1] "Student-1" "Student-2" "Student-3"
\end{verbatim}

\begin{Shaded}
\begin{Highlighting}[]
\FunctionTok{class}\NormalTok{(}\DecValTok{42}\NormalTok{)}
\end{Highlighting}
\end{Shaded}

\begin{verbatim}
## [1] "numeric"
\end{verbatim}

\begin{Shaded}
\begin{Highlighting}[]
\FunctionTok{class}\NormalTok{(}\DecValTok{42}\DataTypeTok{L}\NormalTok{)}
\end{Highlighting}
\end{Shaded}

\begin{verbatim}
## [1] "integer"
\end{verbatim}

\begin{Shaded}
\begin{Highlighting}[]
\FunctionTok{class}\NormalTok{(}\DecValTok{41}\DataTypeTok{i}\NormalTok{)}
\end{Highlighting}
\end{Shaded}

\begin{verbatim}
## [1] "complex"
\end{verbatim}

\begin{Shaded}
\begin{Highlighting}[]
\FunctionTok{class}\NormalTok{(}\ConstantTok{TRUE}\NormalTok{)}
\end{Highlighting}
\end{Shaded}

\begin{verbatim}
## [1] "logical"
\end{verbatim}

\subsubsection{Part B: Data Frames and Conditional
Statements}\label{part-b-data-frames-and-conditional-statements}

\begin{Shaded}
\begin{Highlighting}[]
\CommentTok{\#B1. Create a data frame students with 9 rowsand these columns (use stringsAsFactors = FALSE; sex is "M" or "F".}

\NormalTok{students }\OtherTok{\textless{}{-}} \FunctionTok{data.frame}\NormalTok{(}
  \AttributeTok{id =} \FunctionTok{c}\NormalTok{(}\DecValTok{1}\SpecialCharTok{:}\DecValTok{8}\NormalTok{),}
  \AttributeTok{name =} \FunctionTok{c}\NormalTok{(}\StringTok{"Alice"}\NormalTok{, }\StringTok{"Bruno"}\NormalTok{, }\StringTok{"Chen"}\NormalTok{, }\StringTok{"Dana"}\NormalTok{, }\StringTok{"Elif"}\NormalTok{, }\StringTok{"Farid"}\NormalTok{, }\StringTok{"Grace"}\NormalTok{, }\StringTok{"Hiro"}\NormalTok{),}
  \AttributeTok{sex =} \FunctionTok{c}\NormalTok{(}\StringTok{"F"}\NormalTok{, }\StringTok{"M"}\NormalTok{, }\StringTok{"M"}\NormalTok{, }\StringTok{"F"}\NormalTok{, }\StringTok{"F"}\NormalTok{, }\StringTok{"M"}\NormalTok{, }\StringTok{"F"}\NormalTok{, }\StringTok{"M"}\NormalTok{),}
  \AttributeTok{section =} \FunctionTok{c}\NormalTok{(}\StringTok{"A"}\NormalTok{, }\StringTok{"B"}\NormalTok{, }\StringTok{"A"}\NormalTok{,}\StringTok{"C"}\NormalTok{,}\StringTok{"B"}\NormalTok{,}\StringTok{"C"}\NormalTok{,}\StringTok{"A"}\NormalTok{, }\StringTok{"B"}\NormalTok{),}
  \AttributeTok{midterm =} \FunctionTok{c}\NormalTok{(}\DecValTok{78}\NormalTok{, }\DecValTok{85}\NormalTok{, }\DecValTok{92}\NormalTok{, }\DecValTok{67}\NormalTok{, }\DecValTok{74}\NormalTok{, }\DecValTok{88}\NormalTok{, }\DecValTok{95}\NormalTok{, }\DecValTok{59}\NormalTok{),}
  \AttributeTok{final =} \FunctionTok{c}\NormalTok{(}\DecValTok{82}\NormalTok{, }\DecValTok{79}\NormalTok{, }\DecValTok{96}\NormalTok{, }\DecValTok{71}\NormalTok{, }\DecValTok{80}\NormalTok{, }\DecValTok{91}\NormalTok{, }\DecValTok{93}\NormalTok{, }\DecValTok{65}\NormalTok{),}
  \AttributeTok{attendance =} \FunctionTok{c}\NormalTok{(}\DecValTok{95}\NormalTok{, }\DecValTok{88}\NormalTok{, }\DecValTok{100}\NormalTok{, }\DecValTok{72}\NormalTok{, }\DecValTok{90}\NormalTok{, }\DecValTok{85}\NormalTok{, }\DecValTok{98}\NormalTok{, }\DecValTok{60}\NormalTok{),}
  \AttributeTok{stringsAsFactors =} \ConstantTok{FALSE}
\NormalTok{)}
\NormalTok{students}
\end{Highlighting}
\end{Shaded}

\begin{verbatim}
##   id  name sex section midterm final attendance
## 1  1 Alice   F       A      78    82         95
## 2  2 Bruno   M       B      85    79         88
## 3  3  Chen   M       A      92    96        100
## 4  4  Dana   F       C      67    71         72
## 5  5  Elif   F       B      74    80         90
## 6  6 Farid   M       C      88    91         85
## 7  7 Grace   F       A      95    93         98
## 8  8  Hiro   M       B      59    65         60
\end{verbatim}

\begin{Shaded}
\begin{Highlighting}[]
\CommentTok{\#B2. Display the structure, dimensions, column names, summary of students}

\FunctionTok{str}\NormalTok{(students)}
\end{Highlighting}
\end{Shaded}

\begin{verbatim}
## 'data.frame':    8 obs. of  7 variables:
##  $ id        : int  1 2 3 4 5 6 7 8
##  $ name      : chr  "Alice" "Bruno" "Chen" "Dana" ...
##  $ sex       : chr  "F" "M" "M" "F" ...
##  $ section   : chr  "A" "B" "A" "C" ...
##  $ midterm   : num  78 85 92 67 74 88 95 59
##  $ final     : num  82 79 96 71 80 91 93 65
##  $ attendance: num  95 88 100 72 90 85 98 60
\end{verbatim}

\begin{Shaded}
\begin{Highlighting}[]
\FunctionTok{dim}\NormalTok{(students)}
\end{Highlighting}
\end{Shaded}

\begin{verbatim}
## [1] 8 7
\end{verbatim}

\begin{Shaded}
\begin{Highlighting}[]
\FunctionTok{names}\NormalTok{(students)}
\end{Highlighting}
\end{Shaded}

\begin{verbatim}
## [1] "id"         "name"       "sex"        "section"    "midterm"   
## [6] "final"      "attendance"
\end{verbatim}

\begin{Shaded}
\begin{Highlighting}[]
\FunctionTok{summary}\NormalTok{(students)}
\end{Highlighting}
\end{Shaded}

\begin{verbatim}
##        id              name          sex         section     midterm     
##  Min.   :1.00   Length   :8   Length   :8   Length   :8   Min.   :59.00  
##  1st Qu.:2.75   N.unique :8   N.unique :2   N.unique :3   1st Qu.:72.25  
##  Median :4.50   N.blank  :0   N.blank  :0   N.blank  :0   Median :81.50  
##  Mean   :4.50   Min.nchar:4   Min.nchar:1   Min.nchar:1   Mean   :79.75  
##  3rd Qu.:6.25   Max.nchar:5   Max.nchar:1   Max.nchar:1   3rd Qu.:89.00  
##  Max.   :8.00                                             Max.   :95.00  
##      final         attendance    
##  Min.   :65.00   Min.   : 60.00  
##  1st Qu.:77.00   1st Qu.: 81.75  
##  Median :81.00   Median : 89.00  
##  Mean   :82.12   Mean   : 86.00  
##  3rd Qu.:91.50   3rd Qu.: 95.75  
##  Max.   :96.00   Max.   :100.00
\end{verbatim}

\begin{Shaded}
\begin{Highlighting}[]
\CommentTok{\#B3. Add a column average = the mean of midterm and final for each student, rounded to 1 decimal}

\NormalTok{students}\SpecialCharTok{$}\NormalTok{average }\OtherTok{\textless{}{-}} \FunctionTok{round}\NormalTok{((students}\SpecialCharTok{$}\NormalTok{midterm }\SpecialCharTok{+}\NormalTok{ students}\SpecialCharTok{$}\NormalTok{final) }\SpecialCharTok{/} \DecValTok{2}\NormalTok{, }\DecValTok{1}\NormalTok{)}
\NormalTok{students}
\end{Highlighting}
\end{Shaded}

\begin{verbatim}
##   id  name sex section midterm final attendance average
## 1  1 Alice   F       A      78    82         95    80.0
## 2  2 Bruno   M       B      85    79         88    82.0
## 3  3  Chen   M       A      92    96        100    94.0
## 4  4  Dana   F       C      67    71         72    69.0
## 5  5  Elif   F       B      74    80         90    77.0
## 6  6 Farid   M       C      88    91         85    89.5
## 7  7 Grace   F       A      95    93         98    94.0
## 8  8  Hiro   M       B      59    65         60    62.0
\end{verbatim}

\begin{Shaded}
\begin{Highlighting}[]
\CommentTok{\#B4. Add a logical column passed that is TRUE when average \textgreater{}= 70 and attendance \textgreater{}= 75}

\NormalTok{students}\SpecialCharTok{$}\NormalTok{passed }\OtherTok{\textless{}{-}}\NormalTok{ students}\SpecialCharTok{$}\NormalTok{average }\SpecialCharTok{\textgreater{}=} \DecValTok{70} \SpecialCharTok{\&}\NormalTok{ students}\SpecialCharTok{$}\NormalTok{attendance }\SpecialCharTok{\textgreater{}=} \DecValTok{75}
\NormalTok{students}
\end{Highlighting}
\end{Shaded}

\begin{verbatim}
##   id  name sex section midterm final attendance average passed
## 1  1 Alice   F       A      78    82         95    80.0   TRUE
## 2  2 Bruno   M       B      85    79         88    82.0   TRUE
## 3  3  Chen   M       A      92    96        100    94.0   TRUE
## 4  4  Dana   F       C      67    71         72    69.0  FALSE
## 5  5  Elif   F       B      74    80         90    77.0   TRUE
## 6  6 Farid   M       C      88    91         85    89.5   TRUE
## 7  7 Grace   F       A      95    93         98    94.0   TRUE
## 8  8  Hiro   M       B      59    65         60    62.0  FALSE
\end{verbatim}

\begin{Shaded}
\begin{Highlighting}[]
\CommentTok{\#B5. Using ifelse(), add a column standing with these values based on average.}

\NormalTok{students}\SpecialCharTok{$}\NormalTok{standing }\OtherTok{\textless{}{-}} \FunctionTok{ifelse}\NormalTok{(students}\SpecialCharTok{$}\NormalTok{average }\SpecialCharTok{\textgreater{}=} \DecValTok{90}\NormalTok{, }\StringTok{"Honors"}\NormalTok{, }\FunctionTok{ifelse}\NormalTok{(students}\SpecialCharTok{$}\NormalTok{average }\SpecialCharTok{\textgreater{}=} \DecValTok{75}\NormalTok{, }\StringTok{"Good"}\NormalTok{, }\StringTok{"Needs Support"}\NormalTok{))}
\NormalTok{students}
\end{Highlighting}
\end{Shaded}

\begin{verbatim}
##   id  name sex section midterm final attendance average passed      standing
## 1  1 Alice   F       A      78    82         95    80.0   TRUE          Good
## 2  2 Bruno   M       B      85    79         88    82.0   TRUE          Good
## 3  3  Chen   M       A      92    96        100    94.0   TRUE        Honors
## 4  4  Dana   F       C      67    71         72    69.0  FALSE Needs Support
## 5  5  Elif   F       B      74    80         90    77.0   TRUE          Good
## 6  6 Farid   M       C      88    91         85    89.5   TRUE          Good
## 7  7 Grace   F       A      95    93         98    94.0   TRUE        Honors
## 8  8  Hiro   M       B      59    65         60    62.0  FALSE Needs Support
\end{verbatim}

\begin{Shaded}
\begin{Highlighting}[]
\CommentTok{\#B6. Write a function classify(score) that uses an if/else if/ else chain and returns "Distinction" for 85 and above, "Merit" for 70 to 84, "Pass" for 50 to 69, and "Fail" below 50. Test it with sapply(c(90, 73, 55, 40))}

\NormalTok{classify }\OtherTok{\textless{}{-}} \ControlFlowTok{function}\NormalTok{(score) \{}
  \ControlFlowTok{if}\NormalTok{(score }\SpecialCharTok{\textgreater{}=} \DecValTok{85}\NormalTok{) }\StringTok{"Distinction"}
  \ControlFlowTok{else} \ControlFlowTok{if}\NormalTok{(score }\SpecialCharTok{\textgreater{}=} \DecValTok{70}\NormalTok{) }\StringTok{"Merit"}
  \ControlFlowTok{else} \ControlFlowTok{if}\NormalTok{(score }\SpecialCharTok{\textgreater{}=} \DecValTok{50}\NormalTok{) }\StringTok{"Pass"}
  \ControlFlowTok{else} \StringTok{"Fail"}
\NormalTok{\}}
\FunctionTok{sapply}\NormalTok{(}\FunctionTok{c}\NormalTok{(}\DecValTok{90}\NormalTok{, }\DecValTok{73}\NormalTok{, }\DecValTok{55}\NormalTok{, }\DecValTok{40}\NormalTok{), classify)}
\end{Highlighting}
\end{Shaded}

\begin{verbatim}
## [1] "Distinction" "Merit"       "Pass"        "Fail"
\end{verbatim}

\subsubsection{Part C: rbind() and
cbind()}\label{part-c-rbind-and-cbind}

\begin{Shaded}
\begin{Highlighting}[]
\CommentTok{\#C1. Create a data frame new\_students with two students. Its columns must match the original, including average, passed and standing, computed the same way: id 9, Ines, F, section C, midterm 81, final 84, attendance 92. id 10, Jamal, M, section A, midterm 70, final 68, attendance 80.}
\CommentTok{\# Use rbind() to append them to students, saving the result as students\_all}

\NormalTok{new\_students }\OtherTok{\textless{}{-}} \FunctionTok{data.frame}\NormalTok{(}
  \AttributeTok{id =} \FunctionTok{c}\NormalTok{(}\DecValTok{9}\NormalTok{, }\DecValTok{10}\NormalTok{),}
  \AttributeTok{name =} \FunctionTok{c}\NormalTok{(}\StringTok{"Ines"}\NormalTok{, }\StringTok{"Jamal"}\NormalTok{),}
  \AttributeTok{sex =} \FunctionTok{c}\NormalTok{(}\StringTok{"F"}\NormalTok{, }\StringTok{"M"}\NormalTok{),}
  \AttributeTok{section =} \FunctionTok{c}\NormalTok{(}\StringTok{"C"}\NormalTok{, }\StringTok{"A"}\NormalTok{),}
  \AttributeTok{midterm =} \FunctionTok{c}\NormalTok{(}\DecValTok{81}\NormalTok{, }\DecValTok{70}\NormalTok{),}
  \AttributeTok{final =} \FunctionTok{c}\NormalTok{(}\DecValTok{84}\NormalTok{, }\DecValTok{68}\NormalTok{),}
  \AttributeTok{attendance =} \FunctionTok{c}\NormalTok{(}\DecValTok{92}\NormalTok{, }\DecValTok{80}\NormalTok{),}
  \AttributeTok{stringsAsFactors =} \ConstantTok{FALSE}
\NormalTok{)}

\NormalTok{new\_students}\SpecialCharTok{$}\NormalTok{average }\OtherTok{\textless{}{-}} \FunctionTok{round}\NormalTok{((new\_students}\SpecialCharTok{$}\NormalTok{midterm }\SpecialCharTok{+}\NormalTok{ new\_students}\SpecialCharTok{$}\NormalTok{final) }\SpecialCharTok{/} \DecValTok{2}\NormalTok{, }\DecValTok{1}\NormalTok{)}

\NormalTok{new\_students}\SpecialCharTok{$}\NormalTok{passed }\OtherTok{\textless{}{-}}\NormalTok{ new\_students}\SpecialCharTok{$}\NormalTok{average }\SpecialCharTok{\textgreater{}=} \DecValTok{70} \SpecialCharTok{\&}\NormalTok{ new\_students}\SpecialCharTok{$}\NormalTok{attendance }\SpecialCharTok{\textgreater{}=} \DecValTok{75}

\NormalTok{new\_students}\SpecialCharTok{$}\NormalTok{standing }\OtherTok{\textless{}{-}} \FunctionTok{ifelse}\NormalTok{(new\_students}\SpecialCharTok{$}\NormalTok{average }\SpecialCharTok{\textgreater{}=} \DecValTok{90}\NormalTok{, }\StringTok{"Honors"}\NormalTok{, }\FunctionTok{ifelse}\NormalTok{(new\_students}\SpecialCharTok{$}\NormalTok{average }\SpecialCharTok{\textgreater{}=} \DecValTok{75}\NormalTok{, }\StringTok{"Good"}\NormalTok{, }\StringTok{"Needs Support"}\NormalTok{))}

\NormalTok{students\_all }\OtherTok{\textless{}{-}} \FunctionTok{rbind}\NormalTok{(students, new\_students)}
\NormalTok{students\_all}
\end{Highlighting}
\end{Shaded}

\begin{verbatim}
##    id  name sex section midterm final attendance average passed      standing
## 1   1 Alice   F       A      78    82         95    80.0   TRUE          Good
## 2   2 Bruno   M       B      85    79         88    82.0   TRUE          Good
## 3   3  Chen   M       A      92    96        100    94.0   TRUE        Honors
## 4   4  Dana   F       C      67    71         72    69.0  FALSE Needs Support
## 5   5  Elif   F       B      74    80         90    77.0   TRUE          Good
## 6   6 Farid   M       C      88    91         85    89.5   TRUE          Good
## 7   7 Grace   F       A      95    93         98    94.0   TRUE        Honors
## 8   8  Hiro   M       B      59    65         60    62.0  FALSE Needs Support
## 9   9  Ines   F       C      81    84         92    82.5   TRUE          Good
## 10 10 Jamal   M       A      70    68         80    69.0  FALSE Needs Support
\end{verbatim}

\begin{Shaded}
\begin{Highlighting}[]
\CommentTok{\#C2. Create the numeric vector quiz \textless{}{-} c(8, 7, 10, 6, 7, 9, 10, 5, 8, 6)}

\NormalTok{quiz }\OtherTok{\textless{}{-}} \FunctionTok{c}\NormalTok{(}\DecValTok{8}\NormalTok{, }\DecValTok{7}\NormalTok{, }\DecValTok{10}\NormalTok{, }\DecValTok{6}\NormalTok{, }\DecValTok{7}\NormalTok{, }\DecValTok{9}\NormalTok{, }\DecValTok{10}\NormalTok{, }\DecValTok{5}\NormalTok{, }\DecValTok{8}\NormalTok{, }\DecValTok{6}\NormalTok{)}
\NormalTok{students\_all }\OtherTok{\textless{}{-}} \FunctionTok{cbind}\NormalTok{(students\_all, }\AttributeTok{quiz =}\NormalTok{ quiz)}
\NormalTok{students\_all}
\end{Highlighting}
\end{Shaded}

\begin{verbatim}
##    id  name sex section midterm final attendance average passed      standing
## 1   1 Alice   F       A      78    82         95    80.0   TRUE          Good
## 2   2 Bruno   M       B      85    79         88    82.0   TRUE          Good
## 3   3  Chen   M       A      92    96        100    94.0   TRUE        Honors
## 4   4  Dana   F       C      67    71         72    69.0  FALSE Needs Support
## 5   5  Elif   F       B      74    80         90    77.0   TRUE          Good
## 6   6 Farid   M       C      88    91         85    89.5   TRUE          Good
## 7   7 Grace   F       A      95    93         98    94.0   TRUE        Honors
## 8   8  Hiro   M       B      59    65         60    62.0  FALSE Needs Support
## 9   9  Ines   F       C      81    84         92    82.5   TRUE          Good
## 10 10 Jamal   M       A      70    68         80    69.0  FALSE Needs Support
##    quiz
## 1     8
## 2     7
## 3    10
## 4     6
## 5     7
## 6     9
## 7    10
## 8     5
## 9     8
## 10    6
\end{verbatim}

\begin{Shaded}
\begin{Highlighting}[]
\FunctionTok{class}\NormalTok{(students\_all)}
\end{Highlighting}
\end{Shaded}

\begin{verbatim}
## [1] "data.frame"
\end{verbatim}

\begin{Shaded}
\begin{Highlighting}[]
\CommentTok{\#C3. Create the matrix m \textless{}{-} cbind(midterm = students\_all$midterm, final = students\_all$final) and use rbind() to append a final row containing the column means, labelled "mean".}

\NormalTok{m }\OtherTok{\textless{}{-}} \FunctionTok{cbind}\NormalTok{(}\AttributeTok{midterm =}\NormalTok{ students\_all}\SpecialCharTok{$}\NormalTok{midterm, }\AttributeTok{final =}\NormalTok{ students\_all}\SpecialCharTok{$}\NormalTok{final)}
\NormalTok{m }\OtherTok{\textless{}{-}} \FunctionTok{rbind}\NormalTok{(m, }\AttributeTok{mean =} \FunctionTok{colMeans}\NormalTok{(m))}
\NormalTok{m}
\end{Highlighting}
\end{Shaded}

\begin{verbatim}
##      midterm final
##         78.0  82.0
##         85.0  79.0
##         92.0  96.0
##         67.0  71.0
##         74.0  80.0
##         88.0  91.0
##         95.0  93.0
##         59.0  65.0
##         81.0  84.0
##         70.0  68.0
## mean    78.9  80.9
\end{verbatim}

\begin{Shaded}
\begin{Highlighting}[]
\CommentTok{\#C4. In a comment (2{-}3 sentences), explain what error rbind() produces if the data frames have different column names and show the output if you safely catch it with tryCatch(). tryCatch(rbind(students, bad), error = function(e) \{ message("rbind failed: ", conditionMessage(e));NULL \}).}

\FunctionTok{tryCatch}\NormalTok{(}\FunctionTok{rbind}\NormalTok{(students, bad), }\AttributeTok{error =} \ControlFlowTok{function}\NormalTok{(e) \{ }\FunctionTok{message}\NormalTok{(}\StringTok{"rbind failed: "}\NormalTok{, }\FunctionTok{conditionMessage}\NormalTok{(e));}\ConstantTok{NULL}\NormalTok{ \})}
\end{Highlighting}
\end{Shaded}

\begin{verbatim}
## rbind failed: object 'bad' not found
\end{verbatim}

\begin{verbatim}
## NULL
\end{verbatim}

\begin{Shaded}
\begin{Highlighting}[]
\CommentTok{\# The error rbind() produces if the data frames have different column names is rbind failed: object \textquotesingle{}\textquotesingle{} not found and then it displays null safely instead of stopping the program completely if tryCatch() is used.}
\end{Highlighting}
\end{Shaded}

\subsubsection{Part D: Subsetting}\label{part-d-subsetting}

\begin{Shaded}
\begin{Highlighting}[]
\CommentTok{\# Use students\_all (with quiz) for all questions. Print every result.}
\CommentTok{\#D1. Extract the name column in three ways: $, [["name"]], and [, "name"]}

\NormalTok{students\_all}\SpecialCharTok{$}\NormalTok{name}
\end{Highlighting}
\end{Shaded}

\begin{verbatim}
##  [1] "Alice" "Bruno" "Chen"  "Dana"  "Elif"  "Farid" "Grace" "Hiro"  "Ines" 
## [10] "Jamal"
\end{verbatim}

\begin{Shaded}
\begin{Highlighting}[]
\NormalTok{students\_all[[}\StringTok{"name"}\NormalTok{]]}
\end{Highlighting}
\end{Shaded}

\begin{verbatim}
##  [1] "Alice" "Bruno" "Chen"  "Dana"  "Elif"  "Farid" "Grace" "Hiro"  "Ines" 
## [10] "Jamal"
\end{verbatim}

\begin{Shaded}
\begin{Highlighting}[]
\NormalTok{students\_all[, }\StringTok{"name"}\NormalTok{]}
\end{Highlighting}
\end{Shaded}

\begin{verbatim}
##  [1] "Alice" "Bruno" "Chen"  "Dana"  "Elif"  "Farid" "Grace" "Hiro"  "Ines" 
## [10] "Jamal"
\end{verbatim}

\begin{Shaded}
\begin{Highlighting}[]
\CommentTok{\#D2. Extract rows 2 to 5 and only the columns name, midterm, final.}

\NormalTok{students\_all[}\DecValTok{2}\SpecialCharTok{:}\DecValTok{5}\NormalTok{, }\FunctionTok{c}\NormalTok{(}\StringTok{"name"}\NormalTok{, }\StringTok{"midterm"}\NormalTok{, }\StringTok{"final"}\NormalTok{)]}
\end{Highlighting}
\end{Shaded}

\begin{verbatim}
##    name midterm final
## 2 Bruno      85    79
## 3  Chen      92    96
## 4  Dana      67    71
## 5  Elif      74    80
\end{verbatim}

\begin{Shaded}
\begin{Highlighting}[]
\CommentTok{\#D3. Using logical subsetting, select students with attendance above 90.}

\NormalTok{students\_all[students\_all}\SpecialCharTok{$}\NormalTok{attendance }\SpecialCharTok{\textgreater{}} \DecValTok{90}\NormalTok{, ]}
\end{Highlighting}
\end{Shaded}

\begin{verbatim}
##   id  name sex section midterm final attendance average passed standing quiz
## 1  1 Alice   F       A      78    82         95    80.0   TRUE     Good    8
## 3  3  Chen   M       A      92    96        100    94.0   TRUE   Honors   10
## 7  7 Grace   F       A      95    93         98    94.0   TRUE   Honors   10
## 9  9  Ines   F       C      81    84         92    82.5   TRUE     Good    8
\end{verbatim}

\begin{Shaded}
\begin{Highlighting}[]
\CommentTok{\#D4. Using subset(), select students of sex "F" in section "A" or "C" whose final is at least 75; return only name, section and final}
\FunctionTok{subset}\NormalTok{(students\_all, sex }\SpecialCharTok{==} \StringTok{"F"} \SpecialCharTok{\&}\NormalTok{ section }\SpecialCharTok{\%in\%} \FunctionTok{c}\NormalTok{(}\StringTok{"A"}\NormalTok{, }\StringTok{"B"}\NormalTok{) }\SpecialCharTok{\&}\NormalTok{ final }\SpecialCharTok{\textgreater{}=} \DecValTok{75}\NormalTok{, }\AttributeTok{select =} \FunctionTok{c}\NormalTok{(name, section, final))}
\end{Highlighting}
\end{Shaded}

\begin{verbatim}
##    name section final
## 1 Alice       A    82
## 5  Elif       B    80
## 7 Grace       A    93
\end{verbatim}

\begin{Shaded}
\begin{Highlighting}[]
\CommentTok{\#D5. Use which.max() and which.min() to report the name of the student with the highest and lowest average.}

\NormalTok{highest }\OtherTok{\textless{}{-}}\NormalTok{ students\_all}\SpecialCharTok{$}\NormalTok{name[}\FunctionTok{which.max}\NormalTok{(students\_all}\SpecialCharTok{$}\NormalTok{average)]}
\NormalTok{lowest }\OtherTok{\textless{}{-}}\NormalTok{ students\_all}\SpecialCharTok{$}\NormalTok{name[}\FunctionTok{which.min}\NormalTok{(students\_all}\SpecialCharTok{$}\NormalTok{average)]}
\NormalTok{highest}
\end{Highlighting}
\end{Shaded}

\begin{verbatim}
## [1] "Chen"
\end{verbatim}

\begin{Shaded}
\begin{Highlighting}[]
\NormalTok{lowest}
\end{Highlighting}
\end{Shaded}

\begin{verbatim}
## [1] "Hiro"
\end{verbatim}

\begin{Shaded}
\begin{Highlighting}[]
\CommentTok{\#D6. Sort students\_all by average in descending order, then print the top 3 rows.}

\NormalTok{sort }\OtherTok{\textless{}{-}}\NormalTok{ students\_all[}\FunctionTok{order}\NormalTok{(students\_all}\SpecialCharTok{$}\NormalTok{average, }\AttributeTok{decreasing =} \ConstantTok{TRUE}\NormalTok{), ] }
\FunctionTok{head}\NormalTok{(sort, }\DecValTok{3}\NormalTok{)}
\end{Highlighting}
\end{Shaded}

\begin{verbatim}
##   id  name sex section midterm final attendance average passed standing quiz
## 3  3  Chen   M       A      92    96        100    94.0   TRUE   Honors   10
## 7  7 Grace   F       A      95    93         98    94.0   TRUE   Honors   10
## 6  6 Farid   M       C      88    91         85    89.5   TRUE     Good    9
\end{verbatim}

\subsubsection{Part E: Basic Factors}\label{part-e-basic-factors}

\begin{Shaded}
\begin{Highlighting}[]
\CommentTok{\# Convert section and sex in students\_all to factors with factor(). Print levels() of each.}

\NormalTok{students\_all}\SpecialCharTok{$}\NormalTok{section }\OtherTok{\textless{}{-}} \FunctionTok{factor}\NormalTok{(students\_all}\SpecialCharTok{$}\NormalTok{section)}
\NormalTok{students\_all}\SpecialCharTok{$}\NormalTok{sex }\OtherTok{\textless{}{-}} \FunctionTok{factor}\NormalTok{(students\_all}\SpecialCharTok{$}\NormalTok{sex)}
\FunctionTok{levels}\NormalTok{(students\_all}\SpecialCharTok{$}\NormalTok{section)}
\end{Highlighting}
\end{Shaded}

\begin{verbatim}
## [1] "A" "B" "C"
\end{verbatim}

\begin{Shaded}
\begin{Highlighting}[]
\FunctionTok{levels}\NormalTok{(students\_all}\SpecialCharTok{$}\NormalTok{sex)}
\end{Highlighting}
\end{Shaded}

\begin{verbatim}
## [1] "F" "M"
\end{verbatim}

\begin{Shaded}
\begin{Highlighting}[]
\CommentTok{\#E2. Use table() to count students per sex and per section. Try the two{-}way table of sex by section. table(students\_all$sex, students\_all$section). What is the output?}

\FunctionTok{table}\NormalTok{(students\_all}\SpecialCharTok{$}\NormalTok{sex, students\_all}\SpecialCharTok{$}\NormalTok{section)}
\end{Highlighting}
\end{Shaded}

\begin{verbatim}
##    
##     A B C
##   F 2 1 2
##   M 2 2 1
\end{verbatim}

\begin{Shaded}
\begin{Highlighting}[]
\CommentTok{\#E3. Print as.integer() of the sex factor. In a comment explain what the integers represent.}

\FunctionTok{as.integer}\NormalTok{(students\_all}\SpecialCharTok{$}\NormalTok{sex)}
\end{Highlighting}
\end{Shaded}

\begin{verbatim}
##  [1] 1 2 2 1 1 2 1 2 1 2
\end{verbatim}

\begin{Shaded}
\begin{Highlighting}[]
\CommentTok{\# It just means that it is arranged alphabetically with 1 representing Female and 2 representing Male to classify the two sexes.}
\end{Highlighting}
\end{Shaded}

\begin{Shaded}
\begin{Highlighting}[]
\CommentTok{\#E4. Create a new factor sex\_label from sex using factor()}
\end{Highlighting}
\end{Shaded}


\end{document}
